\documentclass[11pt]{article}
\usepackage[margin=1in]{geometry}
\usepackage{amsmath,amssymb,amsthm,mathtools,array,longtable}
\usepackage{hyperref}

\newtheorem{theorem}{Theorem}[section]
\newtheorem{lemma}[theorem]{Lemma}
\newtheorem{proposition}[theorem]{Proposition}
\newtheorem{corollary}[theorem]{Corollary}
\theoremstyle{definition}
\newtheorem{definition}[theorem]{Definition}
\newtheorem{convention}[theorem]{Convention}
\newtheorem{question}[theorem]{Question}
\theoremstyle{remark}
\newtheorem{remark}[theorem]{Remark}

% ---------------------------------------------------------------------------
% THE PAPER'S NOTATION, CHOSEN ONCE.  Section 2 reconciles it with the sources.
% ---------------------------------------------------------------------------
\newcommand{\Z}{\mathbb{Z}}
\newcommand{\Q}{\mathbb{Q}}
\newcommand{\Lam}{\Lambda}
\newcommand{\Fock}{\mathcal{F}}          % Fock space          (never \F)
\newcommand{\Fld}{F}                     % fraction field      (never \F)
\newcommand{\cube}{\{0,1\}}
\newcommand{\hgt}{\operatorname{ht}}     % one height macro    (not \hgtt/\height)
\newcommand{\Cyl}{\mathcal{C}}           % cylinder sets       (Q150c1 \Cset)
\newcommand{\PH}{\mathcal{P}}            % particle-hole conj  (NOT \mathcal{C})
\newcommand{\one}{\mathbb{1}}
\newcommand{\rev}{\operatorname{rev}}
\newcommand{\cpl}{\operatorname{cpl}}
\newcommand{\supp}{\operatorname{supp}}

\title{Commuting ribbon operators:\\ a classification of the weights\\[4pt]
  \large\normalfont[SKELETON --- 17 September 2026. Statements and provenance only;\\
  proofs are lifted from the files named in each section header.]}
\author{Clio}
\date{17 September 2026}

\begin{document}
\maketitle

\begin{abstract}
For $1\le e<f$ let $R_e^{W}$ be the operator on symmetric functions that adds a connected
$e$-ribbon with a weight $W$ depending on the whole shape of the ribbon.  We determine every
pair $(W,\bar W)$ with $[R_e^{W},R_f^{\bar W}]=0$.  Writing $d=\gcd(e,f)$, the answer is a
\emph{tensor} statement: the commuting pairs are exactly those obtained from a single
rank-$d$ \emph{seed} weight by concatenation, subject to one quadratic condition.  The
content of that condition is order-theoretic: the possible supports of a seed are the closed
sets of a Moore family on $\cube^{\,d-1}$, generated by one local flip rule.  Reading the
classification off that lattice gives the rest.  Its maximal proper elements are
$2\lfloor(d-1)/2\rfloor$ explicit ``cylinders'' of size $2^{\,d-3}$, whence a dichotomy: a
nonzero commuting weight either vanishes nowhere or vanishes on at least three quarters of the
ribbon shapes of its rank, with nothing in between and both extremes attained.  Its
one-element closed sets are the \emph{border-split} words, whence, for every $d\ge3$,
operators that add \emph{one ribbon shape and nothing else} and still commute.  For $d\le2$
the lattice has no proper elements at all, no such operator exists, and the classical
Murnaghan--Nakayama sign is forced.  We also identify the symmetry: particle--hole
conjugation of the fermionic Fock space conjugates $R_e^{W}$ to $R_e^{W\circ T}$, $T$ being
ribbon transposition, and it halves the single-shape classification.
\end{abstract}

\tableofcontents

\bigskip\hrule\medskip
\noindent\textbf{Status of this file.}  Skeleton: it fixes the paper's notation once
(\S\ref{sec:notation}), states the results in full and in their final order
(\S\ref{sec:results}), and records for each later section which artifact the proof is lifted
from, by filename and label (\S\ref{sec:skeleton}).  Nothing here is proved; everything here is
\emph{stated} in the paper's notation, which is not the notation of any one source.
\S\ref{sec:notprove} is the honest list of what remains open.  \S\ref{sec:provenance} records
the four points at which a source is not quoted verbatim; as of 18 September 2026 all four are
\emph{applied} in the text above it, and that section is a reader's note, not a to-do list.
\medskip\hrule\bigskip

\section{Introduction}\label{sec:intro}

\textbf{To be written last.}  The arc: Murnaghan--Nakayama says that multiplying by a power
sum adds ribbons with the sign $(-1)^{\hgt}$, and that these operators commute.  Ask which
\emph{other} weights make them commute.  If the weight may see only the height of the ribbon,
the answer is: none but Murnaghan--Nakayama (\S\ref{sec:height}).  If it may see the whole
shape, the answer is strictly larger as soon as $\gcd(e,f)\ge3$, and the extra solutions are
not deformations of Murnaghan--Nakayama --- they are operators with large zero sets, the
extreme case adding a single ribbon shape (\S\ref{sec:single}).  The organising fact is that
the one-bead sector of the commutator is literally the statement that two words commute in a
free monoid, so the $\gcd$ in the answer is the $\gcd$ of Lyndon--Sch\"utzenberger.

\section{Notation, fixed once}\label{sec:notation}

\begin{convention}[the objects]\label{conv:objects}
$K$ is an integral domain with fraction field $\Fld$.  $\Lam=\Lam_K$ is the ring of symmetric
functions with Schur basis $\{s_\lambda\}$ and Hall inner product.  For a partition $\lambda$,
$M(\lambda)=\{\lambda_j-j:j\ge1\}\subset\Z$ is its \emph{Maya set}; elements are \emph{beads},
elements of $\Z\setminus M$ \emph{holes}.  $\Fock$ is the charge-$0$ semi-infinite wedge with
basis $|M\rangle$, identified with $\Lam$ by $|M(\lambda)\rangle\leftrightarrow s_\lambda$;
the free fermions are undeformed, $\{\psi_a,\psi^*_b\}=\delta_{ab}$, and
$n_a=\psi_a\psi^*_a$.
\end{convention}

\begin{convention}[words, shapes, weights]\label{conv:weights}
A legal \emph{rank-$e$ move} in $M$ is $b\mapsto b+e$ with $b\in M$, $b+e\notin M$; its
\emph{occupancy word} is $u\in\cube^{e-1}$, $u_i=[\,b+i\in M\,]$.  We write $|u|=\sum u_i$ and
$\hgt=|u|$, the number of rows of the ribbon minus one.  A \emph{shape weight of rank $e$} is
a function $W:\cube^{e-1}\to K$; it defines
\[
   R_e^{W}s_\lambda=\sum_{\mu}W(u)\,s_\mu ,
\]
the sum over $\mu\supset\lambda$ with $\mu/\lambda$ a connected $e$-ribbon of occupancy word
$u$.  The \emph{sign-normalised} weight is $\sigma(u):=(-1)^{|u|}W(u)$; Murnaghan--Nakayama is
$\sigma\equiv\text{const}$, i.e.\ $R_e^{W}=\sigma_0\,(p_e\cdot)$.  A \emph{height weight} is
one with $W(u)$ depending only on $|u|$.
\end{convention}

\begin{convention}[the three involutions of $\cube^{n}$]\label{conv:invol}
$(\rev u)_i=u_{n+1-i}$, $(\cpl u)_i=1-u_i$, and $T:=\cpl\rev=\rev\cpl$, so
$(u^{T})_i=1-u_{n+1-i}$.  $T$ is ribbon transposition.  $W^{\vee}:=W\circ T$.
\end{convention}

\begin{convention}[concatenation]\label{conv:otimes}
For shape weights $\sigma_1,\sigma_2$ of ranks $m_1,m_2$,
\[
  (\sigma_1\otimes\sigma_2)(y):=\sigma_1\bigl(y_{[1,m_1-1]}\bigr)\,
       \sigma_2\bigl(y_{[m_1+1,m_1+m_2-1]}\bigr),\qquad y\in\cube^{m_1+m_2-1},
\]
the letter $y_{m_1}$ being ignored; $\otimes$ is associative, and $\sigma^{\otimes k}$ is the
$k$-fold power.
\end{convention}

\begin{convention}[the two quadratic conditions]\label{conv:2B}
A rank-$m$ shape weight $\sigma$ satisfies \textbf{(2B$^{*}_m$)} if for all
$\delta\in[1,m-1]$, $\pi,\varsigma\in\cube^{\delta-1}$, $\pi'\in\cube^{m-\delta-1}$,
\[
   \sigma(\pi\,0\,\pi')\,\sigma(\pi'\,0\,\varsigma)=\sigma(\pi\,1\,\pi')\,\sigma(\pi'\,1\,\varsigma);
\]
and \textbf{(2B$^{**}_m$)} if the same holds with the second factor's prefix decoupled.
(2B$^{**}$) $\Rightarrow$ (2B$^{*}$), strictly from $m\ge4$.
\end{convention}

\begin{convention}[border-split words]\label{conv:BS}
Let $w\in\cube^{n}$.  An integer $p\in[0,n-1]$ is a \emph{border length} of $w$ if
$w_{[1,p]}=w_{[n-p+1,n]}$ (so $p=0$ always is one).  Call $w$ \textbf{border-split} if
\[
   w_{p+1}\ne w_{n-p}\qquad\text{for every border length }p\text{ of }w. \tag{BS}
\]
\end{convention}

\begin{remark}[why not ``non-overlapping'']\label{rem:noname}
\cite{Q150c2} calls (BS) ``non-overlapping''.  That term is already taken: in combinatorics on
words a non-overlapping (= unbordered, bifix-free) word is one with no border length
$p\in[1,n-1]$ at all.  (BS) is strictly weaker --- it permits a border provided the two
continuation letters differ.  The families differ from $n=4$ ($d=5$) on: $8$ border-split
words against $6$ unbordered ones, the difference being $0101$ and $1010$.  The counts
$0,2,4,8,12,28,48,104$ for $d=2,\dots,9$ reported in \cite[\S9]{Q150c2} are the (BS) counts,
and are correct; only the name is hazardous.  \textbf{This paper imports the condition and not
the word.}
\end{remark}

\subsection{Reconciliation with the five sources}\label{sec:reconcile}

The five artifacts do not use one convention.  The table records every symbol whose meaning
changes between them.  A row marked $\blacklozenge$ is a genuine collision --- the same symbol
naming two different objects --- and is the reason this section exists.

\medskip
\noindent\begin{longtable}{@{}p{0.11\textwidth}p{0.20\textwidth}p{0.20\textwidth}p{0.21\textwidth}p{0.20\textwidth}@{}}
\hline
symbol & Q91 / Q92 / Q147 & Q149 & Q150c1 / Q150c2 & this paper\\
\hline\endhead
$\blacklozenge$ \verb|\F| & $\mathcal F$, the Fock space & --- & $\mathbb F$, a \emph{field}
  & \verb|\Fock| $=\mathcal F$ (space), \verb|\Fld| $=\Fld$ (field)\\
$\blacklozenge$ $\mathcal C$ & --- & --- & $\Cyl_{j,b}$, cylinder sets (\verb|\Cset|)
  & $\Cyl_{j,b}$ cylinders; particle--hole conjugation is $\PH$, \emph{never} $\mathcal C$\\
$\blacklozenge$ $\tau$ & Q147: $\tau_N$, the height weight in fermionic coordinates,
  $W_N=(-1)^{N}\tau_N$ & --- & the rank-$d$ \emph{seed} weight of the tensor decomposition
  & $\tau$ = seed only.  Q147's $\tau_N$ is eliminated: we state Q147 in $W$/$\sigma$
  coordinates throughout\\
$\blacklozenge$ $w$ & Q147: $w_N$, a height \emph{weight} & --- & Q150c2: $w\in\cube^{n}$,
  a \emph{word} & $w$ = word only; height weights are $W$ with $W(u)=W_{|u|}$\\
$\blacklozenge$ $\gamma$ & --- & $\gamma:\Z/d\to\Fld^{\times}$, a character
  & Q150c1/c2 proofs: $\gamma\in\cube$, a single \emph{letter}
  & $\gamma$ = character; letters are $b$\\
$\blacklozenge$ $\alpha$ & --- & $\alpha(u)$, the composition (\S1);
  \emph{and} $\alpha\in K^{\times}$, a scalar (\S5) & $\alpha(w)$, the composition
  & $\alpha(u)$ = composition; scalars are $c,\bar c$\\
$\blacklozenge$ $k$ & Q92: $k^{e,f}_{b,c}$, the interference index & ---
  & Q150c2: $k$, the tensor power & $k$ = tensor power; the interference index is
  $\varkappa^{e,f}_{b,c}$\\
$\blacklozenge$ $m$ & $m(x)=[x\in M]$; Q147 $m_e=m(b+e)$ & --- & $m_1,m_2$, \emph{ranks}
  & $m(x)$ = bead indicator; ranks are $e,f,d$\\
$\blacklozenge$ $N$ & $N_{(x,y)}=\sum_{x<j<y}n_j$, an \emph{operator}; Q147 $N$ = a height
  \emph{index} & --- & --- & $N_{(x,y)}$ operator only; height indices are $h$\\
\verb|\Mult{p_e}| & $M_{p_e}$ & $p_e\,\cdot$ & $p_e\,\cdot$ & $p_e\,\cdot$ (macro dropped)\\
height macro & \verb|\hgt| & \verb|\hgtt| & \verb|\height| (c1), \verb|\hgt| (c2)
  & \verb|\hgt|, one macro\\
the operator & $R_e(t)$ (weight $t^{\hgt}$ built in) & $R_e^{W}$ & $R_e^{W}$ & $R_e^{W}$;
  $R_e(t):=R_e^{W}$ with $W(u)=t^{|u|}$\\
$\sigma$ & --- & $\sigma(u)=(-1)^{|u|}W(u)$ & same; $\sigma,\bar\sigma$ of ranks $e,f$
  & same\\
$d$ & --- & $d=\gcd(e,f)$ & $=\gcd(e,f)$ (c1); the rank of a single weight (c2)
  & $d=\gcd(e,f)$ always; a bare rank is $m$\\
$n$ & --- & --- & Q150c2: $n=d-1$, the word length & $n=$ word length, stated each time\\
(NB) & --- & --- & Q150c2 \emph{redefines} ``non-overlapping'' & (BS), border-split
  (Rem.~\ref{rem:noname})\\
\hline
\end{longtable}

\begin{remark}[one sign convention, stated once]\label{rem:signs}
Q147 records $W_N=(-1)^{N}\tau_N$: the $(-1)^{N}$ inside $(-t)^{N}$ in the fermionic normal
form is not part of the weight, it cancels the Jordan--Wigner sign.  We never write $\tau_N$,
so this identity is folded into Convention~\ref{conv:weights} and never reappears.  The one
place the sign is visible is $\sigma=(-1)^{|u|}W$: in $\sigma$-coordinates every matrix element
of the commutator is a difference of two products with $+$ signs.
\end{remark}

\section{The results}\label{sec:results}

The classification (Theorem~\ref{thm:main}) reduces the whole commuting locus to a single
object: one shape weight $\tau$ of rank $d=\gcd(e,f)$, subject to one quadratic condition
(2B$^{*}$).  Everything else --- the dichotomy, the single-shape operators, the counts --- is
a statement about the solutions of that condition, and those solutions are governed by an
order-theoretic structure.  We therefore state the order theory first
(\S\ref{sec:res-order}) and feed it back into the operator statements afterwards
(\S\ref{sec:res-class}).  Throughout, $n:=d-1$ is the word length.

\subsection{The admissible sets form a closure system}\label{sec:res-order}

\begin{definition}[admissible sets]\label{def:adm}
$Z\subseteq\cube^{n}$ is \emph{admissible} (for rank $d=n+1$) if for every $\delta\in[1,d-1]$
and all $\pi,\varsigma\in\cube^{\delta-1}$, $\pi'\in\cube^{d-\delta-1}$,
\[
  \bigl[\,\pi0\pi'\in Z\ \wedge\ \pi'0\varsigma\in Z\,\bigr]
  \iff
  \bigl[\,\pi1\pi'\in Z\ \wedge\ \pi'1\varsigma\in Z\,\bigr].
\]
\end{definition}

\begin{proposition}[admissible $=$ support of a solution]\label{prop:supp}
Let $\sigma$ be a rank-$d$ shape weight over the domain $K$ and $Z=\supp(\sigma)$.  If
$\sigma$ satisfies \textnormal{(2B$^{*}_d$)} then $Z$ is admissible; conversely if $Z$ is
admissible then $\one_Z$ satisfies \textnormal{(2B$^{*}_d$)}.  Hence the admissible subsets of
$\cube^{n}$ are exactly the supports of \textnormal{(2B$^{*}_d$)}-solutions.
\end{proposition}

So the combinatorics of the commuting locus is the combinatorics of the family of admissible
sets.  That family is not merely a list: it is closed under intersection, and it is generated
by a single local rule.

\begin{definition}[linked pairs and the flip]\label{def:linked}
For $u,v\in\cube^{n}$ and $\delta\in[1,n]$, the pair $(u,v)$ is \emph{$\delta$-linked} if
\[
   u_{[\delta+1,n]}=v_{[1,n-\delta]}\qquad\text{and}\qquad u_\delta=v_{n-\delta+1}.
\]
Writing $x^{(i)}$ for $x$ with the letter at position $i$ flipped, set
$\iota_\delta(u,v):=\bigl(u^{(\delta)},v^{(n-\delta+1)}\bigr)$; this is an involution of the
set of $\delta$-linked pairs.
\end{definition}

\begin{theorem}[generation]\label{thm:closure}
$Z\subseteq\cube^{n}$ is admissible if and only if it is closed under the flip rule
\[
   u,v\in Z\ \text{$\delta$-linked}\quad\Longrightarrow\quad
   u^{(\delta)}\in Z\ \text{ and }\ v^{(n-\delta+1)}\in Z
   \qquad(1\le\delta\le n). \tag{R}
\]
\end{theorem}

\begin{proposition}[the Moore family]\label{prop:moore}
The admissible subsets of $\cube^{n}$ are closed under arbitrary intersections, and
$\cube^{n}$ is itself admissible.  They therefore form a closure system: there is a closure
operator $X\mapsto\langle X\rangle$ sending a set of words to the least admissible set
containing it, and the admissible sets are exactly the sets $\langle X\rangle$.
\end{proposition}

This is the result the paper is organised around.  A biconditional condition
(Definition~\ref{def:adm}) has no reason to define a closure system --- intersections of
solution sets of biconditionals are not solution sets in general --- and it is
Theorem~\ref{thm:closure}, which replaces the biconditional by an implication, that makes it
one.  Everything below is obtained by locating specific elements of this lattice: its maximal
proper elements (Corollary~\ref{cor:maximal}), which give the dichotomy, and its atoms
(Corollary~\ref{cor:selfclosed}), which give the single-shape operators.

\begin{definition}[cylinders]\label{def:cyl}
For $1\le j\le d-1$ with $2j\ne d$ and $b\in\cube$ put
$\Cyl_{j,b}:=\{w\in\cube^{n}:w_j=b,\ w_{d-j}=1-b\}$, so $|\Cyl_{j,b}|=2^{\,d-3}$.  Since
$\Cyl_{j,b}=\Cyl_{d-j,1-b}$ there are exactly $2\lfloor(d-1)/2\rfloor$ distinct cylinders,
and none at all for $d\le2$.
\end{definition}

\begin{theorem}[structure of an admissible set]\label{thm:struct}
Let $d\ge2$ and let $Z\subseteq\cube^{n}$ be admissible.  Then $Z=\emptyset$, or
$Z=\cube^{n}$, or $Z\subseteq\Cyl_{j,b}$ for some $j\in[1,d-1]$ with $2j\ne d$ and some
$b\in\cube$.  In the third case, if $j_0$ is the least such $j$, then $Z$ is invariant under
flipping any single coordinate in $\{1,\dots,j_0-1\}\cup\{d-j_0+1,\dots,d-1\}$.
\end{theorem}

\begin{corollary}[the maximal proper admissible sets are exactly the cylinders]\label{cor:maximal}
Let $d\ge3$.  The maximal proper admissible subsets of $\cube^{n}$ are precisely the
$2\lfloor(d-1)/2\rfloor$ distinct cylinders $\Cyl_{j,b}$.  For $d\le2$ there is no proper
nonempty admissible set at all.
\end{corollary}

Counting the maximal elements now gives the density statement, as a corollary rather than a
theorem.

\begin{corollary}[the quarter dichotomy]\label{cor:quarter}
Let $\tau\not\equiv0$ be a rank-$d$ shape weight over $K$ satisfying
\textnormal{(2B$^{*}_d$)}.  Then \emph{either} $\tau$ vanishes nowhere, \emph{or}
\[
   \#\supp(\tau)\;\le\;2^{\,d-3}\;=\;\tfrac14\cdot2^{\,d-1},
\]
i.e.\ $\tau$ vanishes on \emph{at least} three quarters of the $2^{d-1}$ ribbon shapes of rank
$d$.  There is nothing in between.  Both alternatives occur, and the bound $2^{d-3}$ is
\emph{attained}, by $\one_{\Cyl_{j,b}}$ for any admissible index $(j,b)$; so ``at least'' is
sharp and cannot be improved to ``more than''.  For $d\le2$ only the first alternative occurs.
\end{corollary}

\begin{corollary}[the self-closed words]\label{cor:selfclosed}
For $w\in\cube^{n}$: $\langle\{w\}\rangle=\{w\}$ if and only if $w$ is border-split
(Convention~\ref{conv:BS}).  Equivalently, the singletons among the atoms of the closure
system are exactly the border-split words; their number is $2,4,8,12,28,48,104$ for
$d=3,\dots,9$.
\end{corollary}

Indeed rule (R) applied to $u=v=w$ is nontrivial exactly when $(w,w)$ is $\delta$-linked, i.e.\
when $w$ has period $\delta$ and $w_\delta=w_{n-\delta+1}$ --- which is the negation of (BS).
Corollary~\ref{cor:selfclosed} is what makes Theorem~\ref{thm:single} below a statement about
this lattice rather than a separate computation.

\subsection{The operator statements}\label{sec:res-class}

\begin{theorem}[classification]\label{thm:main}
Let $1\le e<f$, $d=\gcd(e,f)$, and let $W,\bar W$ be shape weights of ranks $e,f$ over $K$
with $W\not\equiv0\not\equiv\bar W$.  Then $[R_e^{W},R_f^{\bar W}]=0$ if and only if there are
a nonzero shape weight $\tau$ of rank $d$ over $\Fld$ and $c,\bar c\in\Fld^{\times}$ with
\[
   \sigma=c\,\tau^{\otimes(e/d)},\qquad \bar\sigma=\bar c\,\tau^{\otimes(f/d)},
\]
and $\sigma$ satisfies \textnormal{(2B$^{*}_e$)} (vacuous for $e=1$).
\end{theorem}

\begin{theorem}[single-shape commuting operators]\label{thm:single}
Let $d\ge3$, let $w\in\cube^{d-1}$ be border-split, let $\alpha(w)\models d$ be the
corresponding ribbon shape, let $k\ge2$ and put $e=d$, $f=kd$.  Let $W=\one_w$ and let
$\bar W(y)=(-1)^{\hgt(y)}\prod_{i=1}^{k}\bigl[\,y_{[(i-1)d+1,\,id-1]}=w\,\bigr]$.  Then
$R_d^{W}$ adds the single ribbon shape $\alpha(w)$ and nothing else, $\bar W$ is supported on
exactly $2^{k-1}$ of the $2^{kd-1}$ shapes of rank $kd$, and
$[R_d^{W},R_{kd}^{\bar W}]=0$.  The number of admissible $w$ is
$2,4,8,12,28,48,104$ for $d=3,\dots,9$.
\end{theorem}

By Corollary~\ref{cor:selfclosed} these are exactly the operators whose support is a
\emph{one-element} admissible set; by Corollary~\ref{cor:maximal} the cylinder indicators are
the operators whose support is a \emph{maximal proper} one.  The two extremes of the same
lattice give, respectively, the sparsest commuting operators and the densest vanishing ones.

\begin{theorem}[the symmetry]\label{thm:symmetry}
Let $\PH(M)=\{-1-x:x\in\Z\setminus M\}$ be particle--hole conjugation, an involution of
charge-$0$ Maya sets realising $s_\lambda\mapsto s_{\lambda'}$.  Then for every $e$ and $W$,
\[
   \PH\,R_e^{W}\,\PH\;=\;R_e^{W\circ T},
\]
so $X\mapsto\PH X\PH$ is an algebra automorphism preserving the commuting locus.  On the
nowhere-zero part of Theorem~\ref{thm:main} it acts by a scalar; on the single-shape part it
acts by the involution $w\mapsto w^{T}$ of border-split words, whose fixed points are exactly
the antipalindromes, i.e.\ exactly the $w$ for which $\one_w$ satisfies
\textnormal{(2B$^{**}_d$)}.  Hence the single-shape classification may be carried out on
$T$-orbits: $2,2,6,6,18,24,60$ orbits for $d=3,\dots,9$, against
$2,4,8,12,28,48,104$ words and $2,0,4,0,8,0,16$ antipalindromes.
\end{theorem}

\begin{remark}[a corrected count]\label{rem:orbitfix}
\cite[V10]{Q151} reports the orbit counts as $1,2,2,6,6,18,24,60$ for $d=2,\dots,9$.  The
leading $1$ is wrong and is corrected here to $0$, i.e.\ the list is reindexed to start at
$d=3$.  There are no border-split words of length $n=1$ at all --- the empty border $p=0$ forces
$w_1\ne w_1$ --- so the same source's own shape counts begin $0,2,4,\dots$ at $d=2$, and an
empty set has no orbits.  The remaining seven terms are correct: by Burnside the orbit count is
$\tfrac12(\#\text{words}+\#\text{antipalindromes})$, and $(2{+}2)/2,(4{+}0)/2,(8{+}4)/2,
(12{+}0)/2,(28{+}8)/2,(48{+}0)/2,(104{+}16)/2 = 2,2,6,6,18,24,60$ as stated.  The error is
confined to a value of $d$ at which Theorem~\ref{thm:single} does not apply, so nothing
downstream of it changes.
\end{remark}

\begin{corollary}[the coprime headline]\label{cor:coprime}
If $\gcd(e,f)=1$ then $\tau$ has rank $1$, so $\sigma$ and $\bar\sigma$ are constant:
$R_e^{W}=c\,(p_e\cdot)$ and $R_f^{\bar W}=\bar c\,(p_f\cdot)$.  Shape dependence is forced
away, and with it every parameter but one per rank.
\end{corollary}

\begin{corollary}[height weights]\label{cor:height}
Restricted to height weights the classification collapses for all $e<f$: $[R_e^{W},R_f^{\bar
W}]=0$ with both nonzero forces $W(u)=(-1)^{|u|}W_0$ and $\bar W(y)=(-1)^{|y|}\bar W_0$.
\end{corollary}

\begin{remark}[the threshold, conceptually]\label{rem:threshold}
Definition~\ref{def:cyl} is empty for $d\le2$ and nonempty for $d\ge3$, because
$d\ge3$ is exactly the condition that $[1,d-1]$ contains an index $j$ with $j\ne d-j$.  That
single fact is the whole of the threshold $\gcd(e,f)\ge3$: below it
Theorem~\ref{thm:struct} leaves only $\emptyset$ and $\cube^{n}$, so every commuting weight is
nowhere zero and Corollary~\ref{cor:coprime}'s conclusion propagates; at and above it the
lattice acquires proper elements, and with them the operators of
Theorem~\ref{thm:single}.
\end{remark}

\section{Section skeleton and provenance}\label{sec:skeleton}

Each subsection below states what the corresponding section of the paper proves and the
file and label the proof is lifted from.  Numbers are given from the label map of
Appendix~\ref{app:labels}, generated from the \texttt{.aux} files on 17 September 2026.

\subsection*{\S4 The abacus dictionary and the fermionic normal form}
\emph{Proves:} adding an $e$-ribbon is a bead move $b\mapsto b+e$; the height is the number of
beads jumped; the occupancy word determines the ribbon shape, \emph{top row first}, by
$\alpha(u)=(c_{k+1}-c_k,\dots,c_1-c_0)$ where the $c_i$ list $\{0\}\cup\{i:u_i=1\}\cup\{e\}$;
and $R_e(t)=\sum_b\psi_{b+e}\psi^*_b\,(-t)^{N_{(b,b+e)}}$.
\emph{Lifted from:} \texttt{2026-09-06-c2-Q91-fermionic-normal-form.tex},
\texttt{lem:abacus}, \texttt{lem:sign}, \texttt{lem:dressing},
\texttt{thm:main} (=Thm.~7); and \texttt{2026-09-15-c1-Q149-shape-weights.tex},
\texttt{lem:dict} (=Lem.~1.3).
\emph{Note:} the reversal in $\alpha(u)$ is not cosmetic --- without it the two L-trominoes
exchange names, which is exactly the distinction the cheapest test of the paper turns on.

\subsection*{\S5 The two sectors of the commutator}
\emph{Proves:} $\langle s_\mu,[R_e^{W},R_f^{\bar W}]s_\lambda\rangle=0$ unless
$|M(\lambda)\triangle M(\mu)|\in\{2,4\}$; closed forms for the one-bead and two-bead matrix
elements in $\sigma$-coordinates.
\emph{Lifted from:} \texttt{2026-09-15-c1-Q149-shape-weights.tex}, \texttt{lem:sectors}
(=Lem.~4.1), \texttt{lem:onebead} (=Lem.~4.2), \texttt{lem:twobead} (=Lem.~4.3).
\emph{Background, for the single-$t$ case only:}
\texttt{2026-09-07-Q92-cross-rank-commutator.tex}, \texttt{thm:matrix} (=Thm.~5), giving
$(m(b+e)-m(b+f))(1+t)t^{N-1}$ in the one-bead sector and $t^{P+Q}(t^{-\varkappa}-t^{\varkappa})$
in the two-bead sector, and \texttt{thm:main} (=Thm.~6),
$[R_e(t),R_f(t)]=-\frac{1+t}{t}(\Phi_{e,f}+(t-1)\Psi_{e,f})$.
\emph{Note:} the witness $\langle s_\mu,[R_e(t),R_f(t)]s_\lambda\rangle=t^{f}-t^{f-2}$ is
\texttt{prop:psiwitness} (=Prop.~16) and belongs to the \textbf{two}-bead sector; the brief
for this session filed it under the one-bead sector.

\subsection*{\S6 Height weights force Murnaghan--Nakayama}\label{sec:height}
\emph{Proves:} Corollary~\ref{cor:height}, using only the one-bead sector and only the
occupancy pattern $m_e=1$, $m_f=0$.
\emph{Lifted from:} \texttt{2026-09-11-c2-Q147-multi-t-bracket.tex}, \texttt{thm:main}
(=Thm.~7) via \texttt{lem:B}, with \texttt{lem:realise} as the control that every window
configuration occurs.
\emph{Note:} restate in $W$-coordinates; do not import $\tau_N$
(Remark~\ref{rem:signs}).  The $(1+t)$ factor of Q92 is the $N=0$ instance of the forcing
relation $W_{N+1}=-W_N$.

\subsection*{\S7 The one-bead sector is commutation in a free monoid}
\emph{Proves:} the one-bead sector vanishes identically iff $\bar\sigma=\sigma\otimes\rho
=\rho\otimes\sigma$ for a unique nonzero $\rho$ of rank $f-e$; and two shape weights commute
under $\otimes$ iff both are powers of a common seed --- Lyndon--Sch\"utzenberger for the
semigroup of Convention~\ref{conv:otimes}, by Euclidean descent.  This is where the $\gcd$
comes from.
\emph{Lifted from:} \texttt{2026-09-16-c1-Q150-nonvanishing.tex}, \texttt{prop:1B}
(=Prop.~3.1) and \texttt{thm:structure0} (=Thm.~1.2).

\subsection*{\S8 The two-bead sector is (2B$^{*}$), and the classification}
\emph{Proves:} given the one-bead conclusion, the two-bead sector is exactly
(2B$^{*}_e$) for $\sigma$; hence Theorem~\ref{thm:main}.
\emph{Lifted from:} \texttt{2026-09-16-c1-Q150-nonvanishing.tex}, \texttt{prop:2B} and
\texttt{thm:main0}.
\emph{Note:} \texttt{thm:main0} is a biconditional, proved from the two biconditionals
\texttt{prop:1B}, \texttt{prop:2B} and \cite[\texttt{lem:sectors}]{Q149} alone.  In particular
\textbf{both} directions --- including sufficiency, for which \cite[\texttt{prop:suff}]{Q149}
would be the alternative source --- are available without any appeal to nonvanishing.  This is
why \S\ref{sec:provenance}(P2) is an absence of dependency rather than a repair.
\emph{Note:} \texttt{prop:2B}'s converse direction makes a \emph{second, separate} appeal to
$K$ being a domain, with witness $\rho(\psi_0)$; say so, per
\cite[one-domain-appeal]{memo}.

\subsection*{\S9 The structure of an admissible support, and the dichotomy}
\emph{Proves:} an admissible support is empty, everything, or contained in a cylinder
$\Cyl_{j,b}$ with $2j\ne d$; the maximal proper admissible sets are exactly the
$2\lfloor(d-1)/2\rfloor$ cylinders; hence Theorem~\ref{thm:struct},
Corollary~\ref{cor:maximal} and, by counting them, Corollary~\ref{cor:quarter}.  Also
Remark~\ref{rem:threshold}: $d\ge3$ is exactly the condition that $[1,d-1]$ contains some
$j\ne d-j$, which is the conceptual content of the threshold $\gcd(e,f)\ge3$.
\emph{Lifted from:} \texttt{2026-09-16-c2-Q150-H2-single-shape.tex}, \texttt{def:S}
(=Def.~2.2), \texttt{prop:supp} (=Prop.~2.3), \texttt{thm:struct} (=Thm.~6.1),
\texttt{cor:three} (=Cor.~6.2), \texttt{cor:maximal} (=Cor.~6.3), \texttt{cor:quarter}
(=Cor.~6.4).
\emph{Correction applied:} the source's prose gloss says ``more than three quarters'';
Corollary~\ref{cor:quarter} says ``at least'', which is what its own display asserts and what
the cylinder indicator attains.  See \S\ref{sec:provenance}(P3).

\subsection*{\S10 Single-shape weights}\label{sec:single}
\emph{Proves:} $\one_w$ satisfies (2B$^{*}_d$) iff $w$ is border-split; it satisfies
(2B$^{**}_d$) iff $w$ is an antipalindrome, of which there are $2^{(d-1)/2}$ for $d$ odd and
none for $d$ even; hence Theorem~\ref{thm:single}, with the smallest new case $d=4$, $k=2$,
$w=001$, shape $(1,3)$.
\emph{Lifted from:} \texttt{2026-09-16-c2-Q150-H2-single-shape.tex}, \texttt{def:border}
(=Def.~4.1), \texttt{thm:singleton} (=Thm.~4.2), \texttt{rem:endrows} (=Rem.~4.3),
\texttt{thm:singleton2} (=Thm.~4.5), \texttt{thm:main} (=Thm.~5.1), \texttt{cor:48}
(=Cor.~5.2).
\emph{Correction applied:} the source's ``non-overlapping'' is imported as the condition
(BS) and not as the word; see Convention~\ref{conv:BS}, Remark~\ref{rem:noname} and
\S\ref{sec:provenance}(P1).
\emph{Note:} for even $d$ these operators are new precisely because the (2B$^{**}$) route
produces none of them --- for $d$ even the word length $n=d-1$ is \emph{odd}, and an
antipalindrome $w=w^{T}$ of odd length would need $w_{(n+1)/2}=1-w_{(n+1)/2}$ at its middle
letter.  (The paper's \S\ref{sec:notprove}(H1) states this correctly; an earlier draft of this
note said ``even length'', which is the same fact with the parity of $n$ and of $d$ confused.)

\subsection*{\S11 Generation: the admissible sets form a closure system}
\emph{Proves:} Definition~\ref{def:linked}, Theorem~\ref{thm:closure} and
Proposition~\ref{prop:moore} --- $Z$ is admissible iff closed under the flip rule, so the
admissible sets are the closed sets of a Moore family --- together with
Corollary~\ref{cor:selfclosed}, $\langle\{w\}\rangle=\{w\}$ iff $w$ is border-split.
\emph{Lifted from:} \texttt{2026-09-16-c2-Q150-H2-single-shape.tex}, \texttt{def:linked},
\texttt{thm:closure}, \texttt{prop:moore}, and for Corollary~\ref{cor:selfclosed} the remark
immediately following \texttt{prop:moore} there, read together with \texttt{thm:singleton}.
\emph{Note on placement:} in the source this is the \emph{last} section, stated as an
afterthought to the single-shape classification.  It is the first result here.  The proofs are
unchanged; only their order is.  Proposition~\ref{prop:moore}'s proof is three lines and uses
nothing but Theorem~\ref{thm:closure}, so nothing is lost by promoting it.

\subsection*{\S12 Particle--hole conjugation}
\emph{Proves:} Theorem~\ref{thm:symmetry}; that $\PH$-conjugation is an \emph{automorphism},
not an antiautomorphism, so it is not Terwilliger's $\dagger$; and that no invertible graded
linear map implements $\rev$ or $\cpl$ alone, so $T$ is exactly the part of the word-level
$(\Z/2)^2$ that an operator can see.
\emph{Lifted from:} \texttt{2026-09-17-c1-Q151-particle-hole.tex}, \texttt{thm:A},
\texttt{lem:Ctrans}, \texttt{thm:C}, \texttt{thm:D}, \texttt{cor:parity},
\texttt{prop:norev}, \texttt{cor:accounting}.
\emph{Note:} \texttt{thm:A} also upgrades \cite[\texttt{lem:transpose}]{Q149} (=Lem.~1.6) from
\texttt{computed} to \texttt{proved}: the fact that transposition is reverse-and-complement on
occupancy words, previously checked on $1500$ ribbons, has a three-line proof on the bead
lattice.

\section{Provenance note: four places where a source is not quoted verbatim}\label{sec:provenance}

The corrections below have been \emph{applied} in \S\ref{sec:notation} and
\S\ref{sec:results}; this section exists so that a reader holding the sources can see where
this paper and they part company, and why.  None is a gap in a proof.

\begin{itemize}
\item[\textbf{(P1)}] \textbf{The word ``non-overlapping'' is not imported.}
  \cite[\texttt{def:border}]{Q150c2} uses it for the condition called (BS) in
  Convention~\ref{conv:BS}.  In combinatorics on words a non-overlapping (= unbordered,
  bifix-free) word is one with no border length $p\in[1,n-1]$ at all; (BS) is strictly weaker,
  permitting a border provided the two continuation letters differ.  The families first differ
  at $n=4$ ($d=5$): $8$ border-split words against $6$ unbordered ones, the difference being
  $0101$ and $1010$.  The counts $0,2,4,8,12,28,48,104$ for $d=2,\dots,9$ reported there are
  the (BS) counts and are correct; only the name is hazardous.
  (Remark~\ref{rem:noname}.)

\item[\textbf{(P2)}] \textbf{Nothing is quoted from \cite{Q149} \S5, \S7 or \S8.}  This is the
  one item with correctness at stake, and the resolution is stronger than a repair: the
  dependency is \emph{absent}, not corrected.  \cite[\texttt{lem:nonvanish}]{Q149} (``a nonzero
  commuting weight vanishes nowhere'') is \textbf{false} for $\gcd(e,f)\ge3$
  \cite[\texttt{cor:main0}]{Q150c1}, and \cite{Q149} itself flags it as ``the one gap of the
  paper'' for $e\ge3$, adding ``everything after this point assumes it''.  That warning is
  exact.  It does not reach this paper, because every statement of
  \S\ref{sec:res-class} is taken from \cite[\texttt{thm:main0}]{Q150c1}, which is a
  biconditional proved from \cite[\texttt{prop:1B}]{Q150c1} and
  \cite[\texttt{prop:2B}]{Q150c1} --- both biconditionals --- together with
  \cite[\texttt{lem:sectors}]{Q149}, and \emph{uses no nonvanishing input in either
  direction}.  The citations of \cite{Q149} surviving in \S\ref{sec:skeleton} are
  \texttt{lem:dict}, \texttt{lem:transpose}, \texttt{lem:realise}, \texttt{lem:sectors},
  \texttt{lem:onebead}, \texttt{lem:twobead} --- all in \S\S1,3,4 of that paper, all
  unconditional, none touched by the audit.

  \smallskip
  The audit is retained below because \cite{Q149} remains on the shelf and will be cited by
  later work; it records how far the flaw reaches, which the blanket instruction ``re-check all
  of \cite[\S8]{Q149}'' is too coarse to say.  Two of its five results are unconditional, and
  two of those --- \texttt{lem:flip} and \texttt{prop:e1} --- are in any case superseded by
  \cite[\texttt{prop:1B}]{Q150c1} and reproved by \cite[\texttt{cor:coprime}]{Q150c1}
  respectively, so even a reader who wanted them need not go back.

  \medskip
  \begin{tabular}{@{}llp{0.52\textwidth}@{}}
  \hline
  label & \S & status\\
  \hline
  \texttt{lem:flip} (F0)--(F3) & 8 & \textbf{Unconditional.}  Its proof needs only
    $\sigma\not\equiv0\not\equiv\bar\sigma$ and $K$ a domain --- it picks \emph{some} nonzero
    witness $\sigma(A)$ or $\sigma(D)$, never all of them.  Superseded by
    \cite[\texttt{prop:1B}]{Q150c1}.\\
  \texttt{lem:nonvanish} & 8 & \textbf{False for $\gcd(e,f)\ge3$.}  True exactly for
    $\gcd(e,f)\le2$ \cite[\texttt{thm:gcdle2}]{Q150c1}.  Note this \emph{extends}
    \cite{Q149}'s own proved range, which was $e\le2$: e.g.\ $(e,f)=(4,6)$ has $e>2$ but
    $\gcd=2$.\\
  \texttt{lem:const} & 8 & \textbf{Conditional} --- stated ``granting'' the previous line.
    Valid exactly for $\gcd(e,f)\le2$.\\
  \texttt{lem:char} & 8 & \textbf{Conditional}, likewise.\\
  \texttt{prop:e1} & 8 & \textbf{Unconditional}: a cyclic-shift argument on the one-bead
    relation, never mentioning nonvanishing.  Reproved by
    \cite[\texttt{cor:coprime}]{Q150c1}.\\
  \hline
  \texttt{prop:suff} & 7 & \textbf{Unconditional and still true}: the $\gamma$-family does
    commute.  This is the sufficiency half of \texttt{thm:main} there and survives intact ---
    it is the nowhere-zero part of Theorem~\ref{thm:main}, which
    \cite[\texttt{thm:main0}]{Q150c1} also proves.\\
  \texttt{thm:main}, necessity & 5 & \textbf{Valid exactly for $\gcd(e,f)\le2$.}  Superseded
    by \cite[\texttt{thm:main0}]{Q150c1}, which reaches the same conclusion for $\gcd\le2$ and
    the correct larger one for $\gcd\ge3$.\\
  \texttt{cor:size} & 5 & True as stated about the group of admissible $\gamma$, but its
    \emph{title} (``the size of the locus'') is wrong for $d\ge3$: it sizes the
    $\gamma$-torus, which is then only the nowhere-zero part.\\
  \hline
  \end{tabular}

\item[\textbf{(P3)}] \textbf{``At least three quarters'', not ``more than''.}  The display of
  \cite[\texttt{cor:quarter}]{Q150c2} is $\#\supp(\tau)\le2^{d-3}$, which is correct and
  \emph{attained}: $\one_{\Cyl_{j,b}}$ has support exactly $2^{d-3}$ and satisfies
  (2B$^{*}_d$) (checked for all $j,b$ and $3\le d\le8$).  The prose gloss of that corollary
  therefore overstates its own display.  Corollary~\ref{cor:quarter} states the sharp form and
  records the attaining family.

\item[\textbf{(P4)}] \textbf{The $t^{f}-t^{f-2}$ witness is two-bead, not one-bead.}
  \cite[\texttt{prop:psiwitness}]{Q92} derives it from $t^{P+Q}(t^{-\varkappa}-t^{\varkappa})$
  with $P=0$, $Q=f-1$, $\varkappa=-1$.  The one-bead element of that paper is
  $(m(b+e)-m(b+f))(1+t)t^{N-1}$.  The note in \S\ref{sec:skeleton} (\S5) states this at the
  point of use.
\end{itemize}

\section{Verification of the counts}\label{sec:verify}

Four numerical claims are made above: the border-split counts of Theorem~\ref{thm:single} and
Corollary~\ref{cor:selfclosed}, the antipalindrome and orbit counts of
Theorem~\ref{thm:symmetry}, the cylinder count $2\lfloor(d-1)/2\rfloor$ and size $2^{d-3}$ of
Definition~\ref{def:cyl}, and the family/name discrepancy of Remark~\ref{rem:noname}.  All were
re-enumerated on 18 September 2026 \emph{from Convention~\ref{conv:BS} and
Convention~\ref{conv:invol} as printed in this paper}, by code sharing nothing with the sources'
enumerators.  This is the check that matters for (P1): the sources state these counts under a
different name for the condition, so agreement of the numbers is what certifies that the
condition imported here is the condition they counted.

\begin{center}
\begin{tabular}{@{}lllll@{}}
\hline
$d$ & border-split & antipalindromes & $T$-orbits & unbordered\\
\hline
$2$ & $0$ & $0$ & $0$ & $2$\\
$3$ & $2$ & $2$ & $2$ & $2$\\
$4$ & $4$ & $0$ & $2$ & $4$\\
$5$ & $8$ & $4$ & $6$ & $6$\\
$6$ & $12$ & $0$ & $6$ & $12$\\
$7$ & $28$ & $8$ & $18$ & $20$\\
$8$ & $48$ & $0$ & $24$ & $40$\\
$9$ & $104$ & $16$ & $60$ & $74$\\
\hline
\end{tabular}
\end{center}

Every entry agrees with the sources except the $d=2$ orbit count, corrected in
Remark~\ref{rem:orbitfix}.  The last column is the check behind Remark~\ref{rem:noname}: the
two families first separate at $d=5$, $8$ against $6$, the difference being exactly $0101$ and
$1010$.  The orbit column is additionally constrained by Burnside,
$\#\text{orbits}=\tfrac12(\#\text{words}+\#\text{fixed points})$, which the first two columns
determine independently; it holds for $3\le d\le9$ and fails for the reported $d=2$ value.
That identity is what makes the correction decisive rather than a disagreement between two
enumerations.

\begin{remark}[no external base was found]\label{rem:oeis}
It would be better to check these counts against a sequence this project did not generate.
Two candidates were on hand: OEIS \texttt{A299699} (rim-hook / border-strip tableaux of size
$n$: $1,4,13,50,179,744,2975,\dots$) and \texttt{A296561} (the same refined by shape, Heinz
indexed).  \textbf{Neither is the right base, and no substitute was sought.}  They count
border-strip \emph{tableaux} --- chains of ribbon additions filling a partition --- whereas the
counts above enumerate \emph{binary words of length $d-1$} subject to a border condition; the
objects are different, the index means a different thing in each, and the growth rates differ
($\sim4.5^{n}$ against $\sim2^{n}$).  No alignment or shift brings them together, and looking
for a sequence that did match would produce a base chosen for matching, which is worth nothing.
The counts above therefore rest on the internal Burnside consistency and on independent
re-enumeration, and \emph{not} on an external source.  This is recorded as a limitation, not as
a verification.
\end{remark}

\begin{remark}[locality: the intersection with Murnaghan--Nakayama]\label{rem:local}
There is a natural class of weights that the present classification almost entirely misses, and
saying which is worthwhile because it delimits the result.  Korff's Hecke rim-hook weight
\cite[eq.~(Hecke\_wt)]{Korff}, normalised as he normalises it
($\operatorname{prob}(b)=(t-1)\operatorname{wt}(b)$), is in the coordinates of
Convention~\ref{conv:weights}
\[
   \sigma_{\mathrm{Korff}}(u)\;=\;(t-1)\,t^{\,e-1-|u|}\;=\;(t-1)\,t^{\,e-1}\bigl(t^{-1}\bigr)^{|u|} ,
\]
using $\hgt+\operatorname{wd}-1=e$ to trade the column count for the row count.  So it is a
\emph{height} weight of geometric type $A\,B^{|u|}$: every bead the move jumps contributes the
same factor $B=t^{-1}$, independently of \emph{which} beads they are.  Any weight assembled
multiplicatively from per-letter data has this form.  By Corollary~\ref{cor:height},
$[R_e^{W},R_f^{\bar W}]=0$ with both weights of this form forces $B=1$: \textbf{the weights that
are products over single letters and the weights that commute across ranks meet exactly in
Murnaghan--Nakayama.}  In particular no operator of Theorem~\ref{thm:single} is of this form,
for $\one_w$ is supported on one word and $A\,B^{|u|}$ is supported everywhere or nowhere.

This is the sense in which the product form of Theorem~\ref{thm:main} is \emph{not} a locality
statement, despite looking like one: $\tau^{\otimes(e/d)}$ is a product over blocks of length
$d$ which \emph{discards} every $d$-th letter, and it reduces to a product over single letters
only when $d=1$ --- which is precisely Corollary~\ref{cor:coprime}.

\emph{Source caveat.}  The formula for $\sigma_{\mathrm{Korff}}$ above is verified directly
against \cite{Korff} (deep-read; and the operator form there, \texttt{lem:E}, contains a
typographical inversion in its statement --- his own proof, his definitions of $\psi^{\pm}(t)$,
and his consistency condition $\operatorname{prob}=(t-1)\operatorname{wt}$ all give
$E_{ij}(t)=(1-t^{-1})(-t)^{1-r}$, not the printed $(t^{-1}-1)(-t)^{r-1}$).  The further
identification of the class $A\,B^{|u|}$ with ``realisable by a local six-vertex Boltzmann
weight'' rests on \cite[Rmk.~2.9]{Petrov}, which this project knows only at agent-summary
depth.  \textbf{That identification is therefore not claimed here as proved}; the statement
above is about the class $A\,B^{|u|}$, which is defined internally, and only the name
``local'' is borrowed.
\end{remark}

\section{What is not proved}\label{sec:notprove}

\begin{itemize}
\item[\textbf{(H1)}] \textbf{The exact rank-$d$ form of (2B$^{*}_e$) when $e>d$.}
  Theorem~\ref{thm:main} imposes (2B$^{*}_e$) on $\sigma=c\,\tau^{\otimes(e/d)}$, which is
  what the commutator hands over; but which \emph{seeds} $\tau$ have
  $\tau^{\otimes k}$ satisfying (2B$^{*}_{kd}$) is not known in closed form.  What is known
  \cite[\texttt{rem:gap}]{Q150c1}: (2B$^{**}_d$) $\Rightarrow$ (2B$^{*}_e$) $\Rightarrow$
  (2B$^{*}_d$), and \emph{both} implications are strict for $d\ge4$, so the middle condition
  is genuinely new.  Whether it depends on $k$ is unknown; it does not for $k\in\{2,3\}$ at
  $d\le4$.  Theorem~\ref{thm:single} is the first clean fact about it from the other side: for
  even $d$ \emph{no} single-shape $\tau$ satisfies (2B$^{**}_d$) --- there are no
  antipalindromes of odd length --- yet single-shape $\tau$ at $d=4,6,8$ do give commuting
  pairs.  That is exactly why those examples are new, and it is also why the
  (2B$^{**}$) route of \cite[\texttt{thm:counter}]{Q150c1} cannot reach them.  Note this does not
  weaken any theorem above; it is a question about a \emph{sufficient} condition, not about
  the classification.  Theorem~\ref{thm:symmetry} adds one constraint the answer must satisfy:
  the middle condition is $T$-stable, since (2B$^{*}$) is and $\otimes$ commutes with $T$.

\item[\textbf{(H2a)}] \textbf{The interior of a cylinder is not described.}
  Corollary~\ref{cor:quarter} bounds an admissible support by a cylinder and
  \cite[\texttt{cor:maximal}]{Q150c2} shows the cylinders are the maximal proper ones, but the admissible
  subsets \emph{inside} a cylinder are not classified.  The counts are $1,3,11,31,356$.  The
  obvious induction is \textbf{closed}: there is no uniform maximal admissible set one level
  down --- at $m=4$ the maximal proper admissible subsets of a cylinder have sizes
  $\{4,4,5,5,5,5,6,6,7,9,13\}$, so no single object plays the role the cylinders play one
  level up.  \cite[\texttt{thm:closure}]{Q150c2} gives a generation rule but not a normal
  form.

\item[\textbf{(H2b)}] \textbf{No uniform description of the value lattice $L_Z$.}  Given an
  admissible support $Z$, the group of possible value-vectors is not described uniformly, and
  the sporadic $2$-torsion in it is unexplained.

\item[\textbf{(H2c)}] \textbf{The $T$-fixed locus is known only at the two extremes.}
  Theorem~\ref{thm:symmetry} determines the action of $T$ on the nowhere-zero seeds (a scalar)
  and on the single-shape seeds (the free involution $w\mapsto w^{T}$).  For a seed with zeros
  that is not a single shape, whether $\tau\circ T\in\Fld^{\times}\tau$ is open.

\item[\textbf{(H2d)}] \textbf{Reversal alone.}  \cite[\texttt{prop:norev}]{Q151} rules out
  invertible \emph{degree-preserving} implementations of $\rev$; a degree-mixing one is not
  ruled out.  Since $\rev$ acts nontrivially on the locus (it inverts $\gamma$), an
  implementation would be worth more than $\PH$ is.
\end{itemize}

\appendix
\section{Label $\to$ number map}\label{app:labels}

\textbf{Generated mechanically from the sources' \texttt{.aux} files on 18 September 2026},
not written by hand; regenerate before quoting any printed number, since the counters shift
whenever a source is edited.  Every row is a label this paper mentions.  Q91, Q92 and Q147 use
a flat theorem counter; Q149, Q150c1, Q150c2 and Q151 number by section.  The 34 numbers
carried over from the 17 September table were re-checked against the same \texttt{.aux} files
and all 34 agreed.  \textbf{Cite labels, not numbers:} the numbers below are for a reader
holding the PDFs, and they are the only numbers this paper asserts about its sources.

\medskip
\begin{longtable}{@{}llll@{}}
\hline
file & label & no. & what it is\\
\hline\endhead
Q91 & \texttt{lem:abacus} & 4 & ribbon $=$ bead move\\
Q91 & \texttt{lem:sign} & 5 & the Jordan--Wigner sign\\
Q91 & \texttt{lem:dressing} & 6 & the dressing factor\\
Q92 & \texttt{thm:matrix} & 5 & the two sectors, single $t$\\
Q92 & \texttt{prop:psiwitness} & 16 & the $t^{f}-t^{f-2}$ witness (\textbf{two}-bead, P4)\\
Q92 & \texttt{cor:three} & 19 & ---\\
Q147 & \texttt{lem:realise} & 4 & every window occurs\\
Q147 & \texttt{lem:onebead} & 5 & one-bead element, height form\\
Q147 & \texttt{lem:B} & 8 & the forcing relation\\
Q147 & \texttt{lem:sectors} & 11 & ---\\
Q147 & \texttt{lem:twobead} & 12 & ---\\
Q149 & \texttt{lem:dict} & 1.3 & word $\leftrightarrow$ composition\\
Q149 & \texttt{lem:transpose} & 1.6 & transposition $=$ rev-and-cpl\\
Q149 & \texttt{lem:realise} & 3.1 & every window configuration occurs\\
Q149 & \texttt{lem:sectors} & 4.1 & the two sectors, shape form\\
Q149 & \texttt{lem:onebead} & 4.2 & one-bead element, shape form\\
Q149 & \texttt{lem:twobead} & 4.3 & two-bead element, shape form\\
Q149 & \texttt{cor:size} & 5.4 & size of the $\gamma$-torus (P2)\\
Q149 & \texttt{prop:suff} & 7.1 & sufficiency; unconditional (P2)\\
Q149 & \texttt{lem:flip} & 8.1 & (F0)--(F3); unconditional (P2)\\
Q149 & \texttt{lem:nonvanish} & 8.2 & \textbf{false} for $\gcd\ge3$ (P2)\\
Q149 & \texttt{lem:const} & 8.3 & ---\\
Q149 & \texttt{lem:char} & 8.4 & ---\\
Q149 & \texttt{prop:e1} & 8.5 & ---\\
Q150c1 & \texttt{thm:structure0} & 1.2 & Lyndon--Sch\"utzenberger for $\otimes$\\
Q150c1 & \texttt{thm:main0} & 1.3 & the classification (biconditional)\\
Q150c1 & \texttt{cor:main0} & 1.4 & corrected nonvanishing\\
Q150c1 & \texttt{conv:2Bstar} & 2.1 & (2B$^{*}$), (2B$^{**}$)\\
Q150c1 & \texttt{prop:1B} & 3.1 & one-bead $=$ $\otimes$-commutation\\
Q150c1 & \texttt{prop:2B} & 3.3 & two-bead $=$ (2B$^{*}$)\\
Q150c1 & \texttt{thm:gcdle2} & 6.1 & nonvanishing for $\gcd\le2$\\
Q150c2 & \texttt{def:S} & 2.2 & admissible support\\
Q150c2 & \texttt{def:border} & 4.1 & (BS), there called ``non-overlapping'' (P1)\\
Q150c2 & \texttt{thm:singleton} & 4.2 & $\mathbb{1}_w$ is (2B$^{*}$) iff (BS)\\
Q150c2 & \texttt{rem:endrows} & 4.3 & the end rows\\
Q150c2 & \texttt{thm:singleton2} & 4.5 & (2B$^{**}$) iff antipalindrome\\
Q150c2 & \texttt{cor:48} & 5.2 & the smallest new case, $d=4$\\
Q150c2 & \texttt{cor:three} & 6.2 & the $\gcd\le2$ case, reproved\\
Q150c2 & \texttt{rem:sym} & 6.6 & the word-level $(\mathbb{Z}/2)^2$\\
Q151 & \texttt{lem:Ctrans} & 1.4 & $\mathcal{P}$ on occupancy words\\
Q151 & \texttt{thm:A} & 2.1 & $\mathcal{P}R_e^{W}\mathcal{P}=R_e^{W\circ T}$\\
Q151 & \texttt{thm:C} & 4.4 & not Terwilliger's $\dagger$\\
Q151 & \texttt{thm:D} & 4.5 & action on the locus\\
Q151 & \texttt{cor:parity} & 4.6 & the parity count\\
Q151 & \texttt{prop:norev} & 4.7 & no implementation of $\rev$ alone\\
Q151 & \texttt{cor:accounting} & 4.9 & the $(\mathbb{Z}/2)^2$ accounting\\
\hline
\end{longtable}

\begin{thebibliography}{99}
\bibitem{Q91} Clio, \emph{The fermionic normal form of the ribbon operator},
  \texttt{proofs/2026-09-06-c2-Q91-fermionic-normal-form.tex}, 6 September 2026.
\bibitem{Q92} Clio, \emph{The cross-rank ribbon commutator},
  \texttt{proofs/2026-09-07-Q92-cross-rank-commutator.tex}, 7 September 2026.
\bibitem{Q147} Clio, \emph{Free height weights force Murnaghan--Nakayama},
  \texttt{proofs/2026-09-11-c2-Q147-multi-t-bracket.tex}, 11 September 2026.
\bibitem{Q149} Clio, \emph{Shape weights on ribbons},
  \texttt{proofs/2026-09-15-c1-Q149-shape-weights.tex}, 15 September 2026.
\bibitem{Q150c1} Clio, \emph{Q150: the nonvanishing lemma is false},
  \texttt{proofs/2026-09-16-c1-Q150-nonvanishing.tex}, 16 September 2026.
\bibitem{Q150c2} Clio, \emph{Q150 (H2): single-shape commuting operators},
  \texttt{proofs/2026-09-16-c2-Q150-H2-single-shape.tex}, 16 September 2026.
\bibitem{Q151} Clio, \emph{Particle--hole conjugation is not Terwilliger's $\dagger$},
  \texttt{proofs/2026-09-17-c1-Q151-particle-hole.tex}, 17 September 2026.
\bibitem{Korff} C.~Korff, \emph{Cylindric Hecke characters and Gromov--Witten invariants
  via the asymmetric six-vertex model}, \texttt{arXiv:1906.02565}.
\bibitem{Petrov} L.~Petrov, \emph{(Mittag-Leffler lecture notes)},
  \texttt{arXiv:2609.18502}.  \emph{Known to this project at agent-summary depth only; cited
  once, in Remark~\ref{rem:local}, for a statement explicitly not claimed as proved.}
\bibitem{memo} Clio, working memoranda, \texttt{memory/MEMORY.md}.
\end{thebibliography}

\end{document}
